\documentclass[11pt,a4paper]{article}
\usepackage[T1]{fontenc}
\usepackage{lmodern}
\usepackage[margin=26mm]{geometry}
\usepackage{amsmath,amssymb,amsthm,booktabs,array,longtable}
\usepackage{microtype,multicol}
\usepackage[hidelinks]{hyperref}
\setlength{\parskip}{3pt}
\setlength{\emergencystretch}{2em}
\raggedbottom
\allowdisplaybreaks[1]
\clubpenalty=10000
\widowpenalty=10000
\newtheorem{theorem}{Theorem}[section]
\newtheorem{lemma}[theorem]{Lemma}
\newtheorem{corollary}[theorem]{Corollary}
\theoremstyle{remark}
\newtheorem{remark}[theorem]{Remark}
\newcommand{\Z}{\mathbb Z}
\newcommand{\Q}{\mathbb Q}
\newcommand{\N}{\mathbb N}
\newcommand{\F}{\mathbb F}
\newcommand{\ps}[1]{#1[[x]]}
\newcommand{\Rev}{\operatorname{Rev}}
\newcommand{\rev}{^{\langle-1\rangle}}
\newcommand{\Plus}{\operatorname{Plus}}
\newcommand{\Minus}{\operatorname{Minus}}
\newcommand{\coeff}[2]{[x^{#1}]#2}
\newcommand{\E}{\mathcal E}
\Urlmuskip=0mu plus 1mu
\hypersetup{pdftitle={A397588: Binary parity for a differential power series},pdfauthor={Dev Goyal},pdfsubject={Separate manuscript with AI disclosure and Lean proof map}}
\title{A397588\\[4pt]\large Binary parity for a differential power series}
\author{Dev Goyal}
\date{24 October 2026}
\begin{document}
\maketitle
\begin{center}
\fbox{\begin{minipage}{0.9\linewidth}\small
\textbf{AI disclosure.} ChatGPT (OpenAI) generated the mathematical arguments,
Lean proofs, and this manuscript in response to Dev Goyal's requests.
Dev Goyal is the sole human author. The source sequence and conjecture or
observation are credited to Paul D. Hanna.
\textbf{Independent human review is pending.}
\end{minipage}}
\end{center}
\begin{abstract}
The unique normalized solution of $A=x+(xA^2)\!\,{}^{\prime}$ has an odd coefficient precisely at powers of two. We give the integral construction, the complete parity proof, and the additional identity $A\equiv x\pmod3$. The latter verifies a source assertion already stated without a conjecture label.
\end{abstract}

\section{Source, notation, and scope}

The source is Paul D. Hanna's OEIS entry A397588 \cite{oeis588},
retrieved on 24 October 2026. Its initial sequence index is 1.
We also retain the formal zero constant coefficient at index zero.

All series are formal. The set $\N$ includes zero, and $[x^n]U$ denotes a
coefficient. A normalized series is $x+O(x^2)$; $O(x^r)$ means divisibility
by $x^r$, over any coefficient ring used below. We write $U\circ V=U(V(x))$
only when $V(0)=0$. Compositional inversion, denoted $\Rev(U)$ or $U\rev$,
is distinct from multiplicative inversion of a unit. Congruences are
coefficientwise; ordinary products and compositional iterates are distinct.

\section{Formal algebra used in the proof}

\begin{lemma}[Coefficientwise construction]\label{lem:contraction}
Let $K$ be a commutative ring. Suppose an operator $\Phi$ preserves
normalized series and, for $r\ge2$, satisfies
\[
 G-H=O(x^r)\quad\Longrightarrow\quad
 \Phi(G)-\Phi(H)=O(x^{r+1}).
\]
Then $\Phi$ has exactly one normalized fixed point in $\ps K$.
More generally, a residual equation whose degree-$r$ condition determines
the next unknown coefficient with multiplier one has a unique solution
once its initial coefficients are fixed and its lower-degree conditions
hold identically.
\end{lemma}

\begin{proof}
Begin with $G_0=x$ and set $G_{m+1}=\Phi(G_m)$. Normalization gives
$G_1-G_0=O(x^2)$. Inductively,
$G_{m+1}-G_m=O(x^{m+2})$, so each coefficient eventually stabilizes.
Define $G$ using those stable coefficients. The displayed implication
ensures that each coefficient of $\Phi(G)$ equals the corresponding stable
coefficient, hence $G=\Phi(G)$. If two fixed points agree below degree
$r$, they agree through degree $r$; induction proves uniqueness.

For the residual version, set the next coefficient to the negative of the
known residual term. Its multiplier is one, so the construction uses only
ring operations. No earlier equation changes, and every equation is
eventually imposed. The same coefficient comparison proves uniqueness.
\end{proof}

\begin{lemma}[Binary support]\label{lem:binary}
There is a unique normalized integral series $C$ satisfying $C=x+C^2$.
Its reduction modulo two is
\[
 P(x)=\sum_{k\ge0}x^{2^k}.
\]
More generally, for every integer $d\ge1$, the unique zero-constant
solution of $V+V^2=x^d$ in $\ps{\F_2}$ is
\[
 V=P(x^d)=\sum_{k\ge0}x^{d2^k}.
\]
\end{lemma}

\begin{proof}
The operator $G\mapsto x+G^2$ satisfies Lemma~\ref{lem:contraction},
since $G^2-H^2=(G-H)(G+H)$ and $G+H=O(x)$. In characteristic two,
every series $V$ satisfies $V(x)^2=V(x^2)$. This follows by truncating
above any fixed degree and applying the polynomial Frobenius identity.

Thus $V+V^2=x^d$ is a recursion in which the coefficient at an odd positive
index is fixed immediately, while the coefficient at an even index $n$
depends on the coefficient at $n/2<n$. The constant coefficient is fixed
to zero. This proves uniqueness. The proposed sum is well-defined
coefficientwise because its exponents tend to infinity; adding its square
cancels all terms except $x^d$. Set $d=1$ to identify the reduction of $C$.
\end{proof}

\section{A397588: a differential equation and binary parity}

\begin{theorem}\label{thm:differential}
There is exactly one normalized $A\in\ps{\Z}$ satisfying
\begin{equation}\label{eq:differential}
 A=x+(xA^2)'.
\end{equation}
Its coefficients satisfy, for every $n\ge0$,
\[
 [x^n]A\text{ is odd}\quad\Longleftrightarrow\quad
 n=2^k\text{ for some }k\in\N.
\]
Also $A\equiv x\pmod3$, so $3\mid[x^n]A$ for every $n>1$.
\end{theorem}

\begin{proof}
Over any commutative ring, put $\Phi(G)=x+(xG^2)'$. If $G$ is normalized,
then $xG^2=O(x^3)$ and its derivative is $O(x^2)$, so $\Phi$ is normalized.
For normalized $G,H$ with $G-H=O(x^r)$,
\[
 x(G^2-H^2)=x(G-H)(G+H)=O(x^{r+2}).
\]
Formal differentiation lowers the order by at most one, including in
positive characteristic. Thus $\Phi(G)-\Phi(H)=O(x^{r+1})$.
Lemma~\ref{lem:contraction} constructs the unique solution over every
commutative ring, in particular over $\Z$. Its integer recursion is
\[
 a_0=0,\qquad a_1=1,\qquad
 a_n=(n+1)\sum_{j=1}^{n-1}a_ja_{n-j}\quad(n>1).
\]
This also exhibits integrality directly. Mapping coefficients preserves
the equation, and uniqueness identifies the reduction with the
corresponding canonical finite-ring solution.

In characteristic two the Leibniz rule gives $(A^2)'=2AA'=0$, hence
$(xA^2)'=A^2$. Equation~\eqref{eq:differential} becomes $A+A^2=x$.
Lemma~\ref{lem:binary} identifies $A$ modulo two with $P$, proving the
parity criterion. At $n=0$, both sides of that criterion are false.

In characteristic three, $x$ itself satisfies the equation because
$(x^3)'=3x^2=0$. Uniqueness gives $A=x$ in $\ps{\F_3}$, which proves the
additional divisibility assertion.
\end{proof}

The parity statement is the conjecture in A397588 \cite{oeis588}.
The divisibility by three was already stated there without a conjecture
label. It is included here with a proof and a stronger entire-series
identity, with no claim that it is an additional newly solved problem.


\section{Separate Lean proof and checked scope}

The entry file is \nolinkurl{lean/Entries/A397588.lean}. Its explicit
theorem statements and proof terms use the namespace
\nolinkurl{FinalCollection.A397588}. The complete imported implementations
and their shared mathematical dependencies are included in the collection.
The following local theorem names give the proof map:

\begin{multicols}{2}\raggedright\begin{itemize}\setlength{\itemsep}{2pt}\setlength{\parsep}{0pt}
\item \nolinkurl{differential_existence}.
\item \nolinkurl{differential_definition}.
\item \nolinkurl{differential_source_unique}.
\item \nolinkurl{differential_parity}.
\item \nolinkurl{differential_mod_three}.
\item \nolinkurl{differential_divisible_by_three}.
\end{itemize}\end{multicols}

The main implementation module is \nolinkurl{Round3.DifferentialParity}.
The catalogue records the full original endpoint names and the new entry
theorems. The project pins Lean 4.19.0 and Mathlib revision
\nolinkurl{c44e0c8ee63ca166450922a373c7409c5d26b00b}. The collection
report records source recompilation and transitive axiom checks. Accepted
axioms are propositional extensionality, classical choice, and quotient
soundness. All index conclusions above are unbounded; exact finite
coefficient checks in the package provide separate corroboration.

\paragraph{Limits and status.}
The source conjecture begins at n=1. Divisibility by three was already stated without a conjecture label. Analytic asymptotics are not proved.
The manuscript consolidates completed results; it does not assert global
novelty, exclusive priority, peer review, OEIS acceptance, or a new submission.
Standard algebraic ingredients are not claimed as new discoveries.

\begin{thebibliography}{9}
\raggedright
\bibitem{oeis588}
Paul D. Hanna, \emph{OEIS A397588}, 3 July 2026.
\url{https://oeis.org/A397588}. Retrieved 24 October 2026.
\end{thebibliography}
\end{document}
